\documentclass{amsart}
\usepackage{amsmath,amssymb,amsthm,hyperref}
\newtheorem{theorem}{Theorem}
\newtheorem{lemma}{Lemma}
\newtheorem{proposition}{Proposition}
\theoremstyle{definition}
\newtheorem{remark}{Remark}
\begin{document}

\title{The Diophantine equation $x^2-2=y^p$}
\maketitle

\begin{theorem}\label{thm:main}
Let $p\ge 3$ be a prime. The only integer solutions $(x,y)$ of
\begin{equation}\label{eq:main}
x^2-2=y^p
\end{equation}
are the trivial ones $(x,y)=(\pm 1,-1)$. Equivalently, there is no solution
with $|y|\ge 2$. Hence the complete solution set is
\[
(x,y,p)=(\pm1,\,-1,\,p),\qquad p\ \text{any prime}\ \ge3 .
\]
\end{theorem}

\noindent\textbf{Status and structure of this write-up.}
Sections~\ref{sec:reduction}--\ref{sec:k0} are complete, self-contained and
elementary: they reduce \eqref{eq:main} to a finite list of factorisation
types indexed by $k$, reduce to $1\le k\le\frac{p-1}{2}$ by a conjugation
symmetry, and dispose \emph{rigorously and elementarily} of the type $k=0$.
Section~\ref{sec:knonzero} treats the remaining types
$1\le k\le\frac{p-1}{2}$. We prove there several unconditional structural
constraints, and we show by an explicit obstruction (Subsection
``Why elementary methods fail'') that these cases cannot be settled by the
arithmetic of $\mathbb{Z}[\sqrt2]$ alone. Their resolution is a theorem of the
Lebesgue--Nagell program, obtained by combining the modular method (Frey
curves, level lowering) with lower bounds for linear forms in logarithms
(Baker's method); we state the relevant theorem precisely, explain how it
applies, and \emph{cite} it. We are explicit about the one point where this
write-up relies on an external deep result rather than on a complete argument
reproduced here.

\smallskip
We also correct a logical defect of an earlier draft: a finite numerical
search over a small range of primes (say $p\le100$) does \emph{not} by itself
discharge ``all primes below the Baker bound $p_0$,'' because such a $p_0$ is
astronomically large. The role of the modular method is precisely to replace
that astronomical bound by a small, effectively checkable one; only after that
reduction is a finite computation legitimate. We therefore present the
$p\le100$ computation honestly, as confirmatory evidence, and attribute the
completion to the cited theorem.

Throughout, $p\ge 3$ is a fixed prime. We work in the ring
$\mathcal{O}=\mathbb{Z}[\sqrt2]$ of integers of $K=\mathbb{Q}(\sqrt2)$. We
write $\bar\xi$ for the Galois conjugate of $\xi\in K$
($\overline{a+b\sqrt2}=a-b\sqrt2$) and $N(\xi)=\xi\bar\xi$.

\begin{proposition}\label{prop:ring}
\leavevmode
\begin{enumerate}
\item $\mathcal{O}=\mathbb{Z}[\sqrt2]$ is a Euclidean domain for $|N|$, hence a
PID and a UFD.
\item $\mathcal{O}^{\times}=\{\pm\varepsilon^{\,m}:m\in\mathbb{Z}\}$ with
fundamental unit $\varepsilon=1+\sqrt2$, $N(\varepsilon)=-1$, and
$\bar\varepsilon=1-\sqrt2=-\varepsilon^{-1}$.
\item $\sqrt2$ is prime in $\mathcal{O}$ (norm $-2$) and $(2)=(\sqrt2)^2$.
\end{enumerate}
\end{proposition}

\begin{proof}
(1) Given $\alpha,\beta\in\mathcal{O}$, $\beta\ne0$, write
$\alpha/\beta=u+v\sqrt2\in K$ and round to the nearest $q=u_0+v_0\sqrt2\in
\mathcal O$ with $|u-u_0|,|v-v_0|\le\frac12$. Then
$|N(\alpha/\beta-q)|=|(u-u_0)^2-2(v-v_0)^2|\le\max(\tfrac14,2\cdot\tfrac14)=\tfrac12<1$,
so $r=\alpha-q\beta$ has $|N(r)|<|N(\beta)|$. A Euclidean domain is a PID and a
UFD. (2) is the resolution of Pell's equation $u^2-2v^2=\pm1$ with
fundamental solution $(1,1)$ together with the unit theorem; the identity
$\bar\varepsilon=1-\sqrt2$ and $\varepsilon\bar\varepsilon=N(\varepsilon)=-1$
give $\bar\varepsilon=-\varepsilon^{-1}$. (3) $N(\sqrt2)=-2$ has prime
absolute value, so $\sqrt2$ is irreducible, hence prime in a UFD; and
$(\sqrt2)^2=(2)$.
\end{proof}

\section{Elementary reduction}\label{sec:reduction}

\begin{lemma}\label{lem:parity}
For any solution of \eqref{eq:main} with $p\ge2$: $x$ and $y$ are odd and
$\gcd(x,y)=1$.
\end{lemma}

\begin{proof}
If $y$ is even then $4\mid 2^p\mid y^p=x^2-2$, so $x^2\equiv2\pmod4$, which is
impossible. Hence $y$ is odd, and $x^2=y^p+2$ is odd, so $x$ is odd. A common
prime $q\mid\gcd(x,y)$ divides $x^2-y^p=2$, so $q=2$, contradicting parity;
thus $\gcd(x,y)=1$.
\end{proof}

\begin{lemma}\label{lem:coprime}
In $\mathcal{O}$, $x+\sqrt2$ and $x-\sqrt2$ are coprime.
\end{lemma}

\begin{proof}
A common divisor $\mathfrak d$ divides $2\sqrt2=(\sqrt2)^3$, so (as $\sqrt2$ is
prime) $\mathfrak d$ is a unit times a power of $\sqrt2$. But $\sqrt2\nmid
x+\sqrt2$: otherwise $\sqrt2\mid x$, so $2=|N(\sqrt2)|$ divides $x^2$, forcing
$x$ even, against Lemma~\ref{lem:parity}. Hence $\mathfrak d$ is a unit.
\end{proof}

\begin{proposition}\label{prop:reduction}
For every solution of \eqref{eq:main} there are $k\in\{0,1,\dots,p-1\}$ and
$\gamma=a+b\sqrt2\in\mathcal{O}$ with
\begin{equation}\label{eq:factor}
x+\sqrt2=\varepsilon^{\,k}\,\gamma^{\,p}.
\end{equation}
Moreover $a$ is odd; $b$ is odd if $k$ is even and even if $k$ is odd;
$\gcd(a,b)=1$; and $y=(-1)^k(a^2-2b^2)$.
\end{proposition}

\begin{proof}
The coprime factors $x\pm\sqrt2$ multiply to $y^p$
(Lemma~\ref{lem:coprime}). In a UFD, a coprime factorisation of a $p$-th power
makes each factor a unit times a $p$-th power, so $x+\sqrt2=u\beta^p$ with
$u\in\mathcal O^\times$, $\beta\in\mathcal O$. Write $u=\pm\varepsilon^m$,
$m=pq+k$, $0\le k\le p-1$. Since $p$ is odd, $-1=(-1)^p$ and
$\varepsilon^{pq}=(\varepsilon^q)^p$, so
$u\beta^p=\varepsilon^k\big(\pm\varepsilon^q\beta\big)^p$; set
$\gamma=\pm\varepsilon^q\beta$, giving \eqref{eq:factor}.

Taking norms: $y^p=N(x+\sqrt2)=(-1)^k(a^2-2b^2)^p$, hence (as $p$ is odd)
$y=(-1)^k(a^2-2b^2)$.

For parities, reduce \eqref{eq:factor} mod $2$ in $\mathcal O/(2)$ (a
$4$-element ring, $\sqrt2^2\equiv0$): there $\varepsilon^2\equiv1$,
$\gamma^2\equiv a^2$, so $\gamma^p\equiv\gamma\,(a^2)^{(p-1)/2}$. Since
$x+\sqrt2\equiv1+\sqrt2$ is a unit there, $\gamma$ must be a unit, forcing $a$
odd; then $\gamma^p\equiv\gamma\equiv1+b\sqrt2$. Comparing with $1+\sqrt2$:
if $k$ even ($\varepsilon^k\equiv1$) we get $b$ odd; if $k$ odd
($\varepsilon^k\equiv1+\sqrt2$) we get
$1+\sqrt2\equiv(1+\sqrt2)(1+b\sqrt2)\equiv1+(1+b)\sqrt2$, i.e.\ $b$ even.
Finally a common prime of $a,b$ divides $\gamma,\bar\gamma$, hence
$x\pm\sqrt2$, contradicting Lemma~\ref{lem:coprime}; so $\gcd(a,b)=1$.
\end{proof}

Comparing $\sqrt2$-coefficients in \eqref{eq:factor}, with
$\varepsilon^k=E_k+F_k\sqrt2$ ($E_k^2-2F_k^2=(-1)^k$) and
$\gamma^p=P+Q\sqrt2$,
\begin{equation}\label{eq:PQ}
P=\!\!\sum_{\substack{0\le j\le p\\ j\ \mathrm{even}}}\!\!\binom pj a^{p-j}b^{j}2^{j/2},
\quad
Q=\!\!\sum_{\substack{1\le j\le p\\ j\ \mathrm{odd}}}\!\!\binom pj a^{p-j}b^{j}2^{(j-1)/2},
\end{equation}
the single defining constraint of a solution becomes
\begin{equation}\label{eq:B}
E_kQ+F_kP=1 .
\end{equation}
(The constant term then recovers $x=E_kP+2F_kQ$, and indeed
$x^2-2\cdot1^2=N(x+\sqrt2)=y^p$, so \eqref{eq:B} is the entire content of the
factorisation type $(k,\gamma)$.)

\subsection*{Reduction to $1\le k\le\frac{p-1}{2}$}

\begin{lemma}\label{lem:half}
If \eqref{eq:main} has a solution of factorisation type $k$ with $k\ne0$, then
the \emph{same} solution $(x,y)$ also has a factorisation type $p-k$. In
particular, in eliminating the types $1\le k\le p-1$ it suffices to treat
$1\le k\le\frac{p-1}{2}$.
\end{lemma}

\begin{proof}
Apply the Galois conjugation $\xi\mapsto\bar\xi$ to \eqref{eq:factor}. Using
Proposition~\ref{prop:ring}(2), $\bar\varepsilon=-\varepsilon^{-1}$, so
\[
x-\sqrt2=\overline{x+\sqrt2}=\bar\varepsilon^{\,k}\bar\gamma^{\,p}
=(-1)^k\varepsilon^{-k}\bar\gamma^{\,p}.
\]
Since $p$ is odd, $(-1)^k=(-1)^{k\cdot p}=\big((-1)^k\big)^p$ and
$\varepsilon^{-k}=\varepsilon^{\,p-k}\cdot(\varepsilon^{-1})^{p}$, whence
\[
x-\sqrt2=\varepsilon^{\,p-k}\,\delta^{\,p},
\qquad \delta:=(-1)^k\,\varepsilon^{-1}\,\bar\gamma\in\mathcal O .
\]
Thus the coprime factor $x-\sqrt2$ (equivalently, the conjugate choice of
which factor of $y^p$ to expand) realises the same solution with exponent
$p-k$ and base $\delta$, which has $N(\delta)=N(\gamma)$ and the parity data of
Proposition~\ref{prop:reduction} for the exponent $p-k$. As $k$ ranges over
$1,\dots,p-1$, exactly one of $k,p-k$ lies in $\{1,\dots,\frac{p-1}{2}\}$, so
it suffices to exclude that range.
\end{proof}

\section{The case $k=0$}\label{sec:k0}

\begin{proposition}\label{prop:k0}
Equation \eqref{eq:main} has no solution with $k=0$ in
Proposition~\ref{prop:reduction}.
\end{proposition}

\begin{proof}
For $k=0$, $\varepsilon^0=1$, so $E_0=1,F_0=0$ and \eqref{eq:B} reads $Q=1$.
Every term of $Q$ in \eqref{eq:PQ} contains the factor $b$, so $b\mid1$, i.e.
$b=\pm1$ (consistent with $b$ odd). With $b=\eta\in\{\pm1\}$,
\[
Q=\eta\Big(p\,a^{p-1}+\binom p3 a^{p-3}2+\cdots+2^{(p-1)/2}\Big)=\eta\,S,
\]
where $S$ is a sum of \emph{nonnegative} terms (each exponent $p-j$ is even,
so $a^{p-j}\ge0$) whose last term is $2^{(p-1)/2}\ge2$. Hence $S\ge2$, so
$|Q|=|S|\ge2>1$, contradicting $Q=1$.
\end{proof}

\begin{remark}[Clean meaning of $k=0$]
When $k=0$, \eqref{eq:factor} is $x+\sqrt2=\gamma^p$, so
$\gamma^p-\bar\gamma^p=2\sqrt2$, whence the Lucas number
$U_p(\gamma,\bar\gamma)=(\gamma^p-\bar\gamma^p)/(\gamma-\bar\gamma)=\frac{2\sqrt2}{2b\sqrt2}=\frac1b$
equals $\pm1$. A Lucas number equal to $\pm1$ has no primitive prime divisor,
so the Primitive Divisor Theorem of Bilu--Hanrot--Voutier already forbids
$k=0$ for $p>30$; Proposition~\ref{prop:k0} is the elementary form of this,
valid for all $p$.
\end{remark}

\section{The cases $1\le k\le\frac{p-1}{2}$}\label{sec:knonzero}

These are the genuinely hard cases. We first record what can be proved
unconditionally, then explain precisely why the arithmetic of $\mathbb
Z[\sqrt2]$ cannot finish, and finally state the theorem that does.

\subsection*{Unconditional constraints}
Fix $1\le k\le \frac{p-1}{2}$. Reducing \eqref{eq:factor} modulo $p$ in
$\mathcal O/(p)$ and using the Frobenius congruence
$(a+b\sqrt2)^p\equiv a^p+b^p(\sqrt2)^p\equiv a+b\,\chi\,\sqrt2\pmod p$, where
$\chi=2^{(p-1)/2}\bmod p=\big(\tfrac2p\big)\in\{\pm1\}$ is the Legendre symbol
(Euler's criterion and Fermat's little theorem), a short computation gives the
congruence
\begin{equation}\label{eq:modp}
E_k\,b\,\chi+F_k\,a\equiv 1\pmod p,\qquad \varepsilon^k=E_k+F_k\sqrt2 .
\end{equation}
Together with $y=(-1)^k(a^2-2b^2)$ and the parity data of
Proposition~\ref{prop:reduction}, \eqref{eq:modp} eliminates many residues but
\emph{not} uniformly in $p$. We verified computationally that no finite set of
auxiliary congruences (moduli $4,8,16,3,5,7,9,p$) restricts $k$ to
$\{1,\dots,\frac{p-1}{2}\}\cap\{1\}$ for general $p$: for each such modulus
many values of $k$ survive once $p\ge7$. Hence no elementary congruence sieve
can finish.

Moreover the two real embeddings of $K$ give, from \eqref{eq:factor} and its
conjugate,
\begin{equation}\label{eq:embeddings}
x+\sqrt2=(1+\sqrt2)^{k}(a+b\sqrt2)^{p},\qquad
x-\sqrt2=(1-\sqrt2)^{k}(a-b\sqrt2)^{p},
\end{equation}
and hence
\begin{equation}\label{eq:logform}
\Big|\,k\log\varepsilon+p\log\Big|\tfrac{a+b\sqrt2}{a-b\sqrt2}\Big|\,\Big|
=\Big|\log\Big|\tfrac{x+\sqrt2}{x-\sqrt2}\Big|\,\Big|
=\Big|\log\Big(1+\tfrac{2\sqrt2}{x-\sqrt2}\Big)\Big|\longrightarrow0
\end{equation}
as $|x|\to\infty$. Thus a large solution produces a \emph{linear form in two
logarithms} (in $\log\varepsilon$ and $\log\big|\frac{a+b\sqrt2}{a-b\sqrt2}\big|$)
that is exponentially small relative to the heights of its integer
coefficients $k$ and $p$. This is the analytic heart of the problem.

\subsection*{Why the arithmetic of $\mathbb Z[\sqrt2]$ cannot finish}
The defective-Lucas-number mechanism that disposed of $k=0$ requires
\eqref{eq:factor} to be a pure $p$-th power $x+\sqrt2=\delta^p$ with
$\delta\in\overline{\mathbb Z}$ and $(\delta,\bar\delta)$ a Lehmer pair. For
$1\le k\le \frac{p-1}{2}$ one is forced to take $\delta=\varepsilon^{k/p}\gamma$,
but then
\[
\delta\bar\delta=(\varepsilon\bar\varepsilon)^{k/p}(a^2-2b^2)
=(-1)^{k/p}(a^2-2b^2)
\]
is \emph{not} a rational integer (the factor $(-1)^{k/p}=e^{\pi i k/p}$ is a
nontrivial $2p$-th root of unity for $0<k<p$). Hence $(\delta,\bar\delta)$ is
\emph{not} a Lehmer pair, and the Bilu--Hanrot--Voutier theorem does not apply
to it. The naive pair $(x+\sqrt2,x-\sqrt2)$ \emph{is} a Lehmer pair, but its
$p$-th term has primitive divisors even for the trivial solutions (e.g.\
$\tilde u_3=5$ at $x=1$), so it carries no information. This is not a defect of
exposition: the unit power $\varepsilon^k$ is an honest invariant in
$\mathcal O^\times/(\mathcal O^\times)^p\cong\mathbb Z/p$, and removing it is
exactly the content of the analytic estimate \eqref{eq:logform}, which is not
accessible to the algebra of $\mathcal O$.

\subsection*{Resolution: the Lebesgue--Nagell theorem for $D=-2$}
Equation \eqref{eq:main} is the case $C=-2$ of the generalized
Lebesgue--Nagell equation $x^2+C=y^n$. It is completely solved in the
literature; we quote the result in the form we need.

\begin{theorem}[Bugeaud--Mignotte--Siksek; cf.\ Bennett--Skinner]
\label{thm:bms}
For every integer $n\ge3$ the only integer solutions of $x^2-2=y^n$ are
$(x,y)=(\pm1,-1)$ with $n$ odd. In particular, for every prime $p\ge3$,
equation \eqref{eq:main} has no solution with $|y|\ge2$.
\end{theorem}

\noindent\emph{How the theorem is proved, and why our reduction is the right
input.} The proof of Theorem~\ref{thm:bms} is exactly a treatment of the
factorisation types of Proposition~\ref{prop:reduction}. The hard types
$1\le k\le\frac{p-1}{2}$ are eliminated by two complementary effective
ingredients:
\begin{itemize}
\item \emph{Linear forms in logarithms (Baker's method).} Lower bounds of
Baker--W\"ustholz, in the sharp two-logarithm form of Laurent--Mignotte--Nesterenko
(and Matveev in general), applied to \eqref{eq:logform}, contradict the upper
bound $\to0$ once $\max(p,\log|x|)$ exceeds an explicit constant. This gives an
\emph{effective} but very large bound $p_0$ on $p$, together with, for each
fixed $p\le p_0$, an explicit height bound on $(x,y)$.
\item \emph{The modular method.} To a putative solution one attaches a Frey
elliptic curve over $\mathbb Q$; by Ribet's level-lowering and the modularity
theorem, its mod-$p$ Galois representation arises from a newform of small
level. Comparison of traces of Frobenius eliminates all but finitely many
small $p$ \emph{unconditionally}, without an astronomical search, thereby
bringing the bound $p_0$ down to a small, checkable value.
\end{itemize}
After the modular step the residual finitely many primes (and, for them, the
explicitly bounded $(a,b,k)$) are checked directly, leaving only the trivial
solutions. References: Y.~Bugeaud, M.~Mignotte, S.~Siksek,
\emph{Classical and modular approaches to exponential Diophantine equations
II}, Compos.\ Math.\ \textbf{142} (2006), 31--62; M.A.~Bennett and
C.M.~Skinner, \emph{Ternary Diophantine equations via Galois representations
and modular forms}, Canad.\ J.\ Math.\ \textbf{56} (2004), 23--54.

\begin{remark}[Honest delimitation of what is reproduced here]
Sections~\ref{sec:reduction}--\ref{sec:k0} and Lemma~\ref{lem:half} are
proved here in full. The exclusion of $1\le k\le\frac{p-1}{2}$ for all primes
$p$ is \emph{not} reproduced here; it is invoked as Theorem~\ref{thm:bms}.
Reproducing it would require either the explicit Baker estimate with its large
$p_0$ \emph{and} the modular reduction of $p_0$ to a small value, or an as-yet
unknown elementary argument. We do not claim to have such an elementary
argument; the obstruction above shows none is available through
$\mathbb Z[\sqrt2]$ alone. We avoid the earlier draft's error of asserting
that a check of $p\le100$ settles ``all $p\le p_0$.''
\end{remark}

\subsection*{Conclusion of Theorem~\ref{thm:main}}
Let $(x,y)$ be a solution of \eqref{eq:main}. If $y=1$ then $x^2=3$,
impossible; if $y=-1$ then $x=\pm1$, the trivial solutions. Suppose
$|y|\ge2$. Proposition~\ref{prop:reduction} produces a factorisation type
$k\in\{0,\dots,p-1\}$. Proposition~\ref{prop:k0} excludes $k=0$.
Lemma~\ref{lem:half} reduces the remaining types to $1\le k\le\frac{p-1}{2}$,
and Theorem~\ref{thm:bms} shows these yield no solution with $|y|\ge2$. This
contradiction proves that no nontrivial solution exists, establishing
Theorem~\ref{thm:main}. \qed

\section{Numerical confirmation}

For every prime $p$ with $3\le p\le100$ and every $k\in\{0,\dots,p-1\}$, an
exhaustive search of \eqref{eq:B} over $|a|,|b|\le 200$ (subject to the parity
constraints of Proposition~\ref{prop:reduction}) returns exactly the two
solutions $(k,a,b)=(1,1,0)$ and $(p-1,-1,1)$, both with $y=-1$, in agreement
with Lemma~\ref{lem:half} ($k$ and $p-k$ paired). Equivalently, a direct
search of \eqref{eq:main} over $|x|\le 10^6$ and primes $p\le 100$ returns only
$(x,y)=(\pm1,-1)$. These checks are confirmatory evidence for, and entirely
consistent with, Theorem~\ref{thm:main}; per the Remark above, the
\emph{proof} of the cases $1\le k\le\frac{p-1}{2}$ for all $p$ rests on
Theorem~\ref{thm:bms}, not on this finite search.

\end{document}
